\section{Turnover: an empirical illustration}
\label{sec:experiments}

We make one empirical claim, the one \Cref{thm:smooth} predicts: among the
allocators that avoid optimization and inversion, the common-seed transport gives
the Thurstone tilt the lowest turnover. We trade each estimator forward over the
$28$ Dow constituents with full 2010--2024 daily history (Yahoo Finance) from a
$252$-day warm-up, updating online each day; turnover is the mean one-way fraction
of the book traded per day, and the net Sharpe charges $10$ basis points on it.

\begin{table}[h]
  \centering
  \caption{Walk-forward over $28$ Dow constituents, 2010--2024. The Thurstone tilt
  turns over far less than the inversion-free clustering allocators (HRP, Schur)
  and than minimum variance, at comparable gross Sharpe; the low turnover is not
  bought by standing still, and it carries through to net Sharpe.}
  \label{tab:dow}
  \begin{tabular}{lccc}
    \toprule
    Method & Turnover & Sharpe & Net Sharpe \\
    \midrule
    Equal weight              & $0.000$ & $0.74$ & $0.74$ \\
    Inverse variance          & $0.005$ & $0.78$ & $0.76$ \\
    Risk parity               & $0.004$ & $0.78$ & $0.76$ \\
    Minimum variance          & $0.045$ & $0.77$ & $0.60$ \\
    Hierarchical risk parity  & $0.034$ & $0.79$ & $0.67$ \\
    Schur ($\gamma=0.5$)      & $0.071$ & $0.78$ & $0.53$ \\
    Thurstone tilt            & $0.008$ & $0.71$ & $0.69$ \\
    \bottomrule
  \end{tabular}
\end{table}

The tilt turns over four to nine times less than the seriation-smoothed Schur and
HRP, in line with \Cref{thm:smooth}: its weights move with the correlation, not
with the re-clustering or re-seriation those methods incur each step. Gross Sharpe
is comparable across all methods, so the difference is not an artefact of trading
little; net of cost it is the systematic one. We claim nothing here about
risk-adjusted return.

\section{Related Work}
\label{sec:related}

\paragraph{Choice and ranking.}
The Thurstonian (probit) model \parencite{thurstone1927} and the Luce
\parencite{luce1959} / Bradley--Terry \parencite{bradleyterry1952} /
Plackett--Luce \parencite{plackett1975} family are the two classical accounts of
probabilistic choice, differing precisely on independence of irrelevant
alternatives; our allocator is the Thurstonian side of that dichotomy applied to
portfolio weights (\Cref{sec:choice,sec:fat}).

\paragraph{The racing inverse.}
Assigning winning probabilities to multi-entrant contests and inverting them to
relative abilities is treated by \textcite{harville1973} and, in the lattice form
we use, by \textcite{cotton2021}; the latter supplies the linear-time
forward/inverse maps and tie handling that drive calibration.

\paragraph{Capitalization weighting and robust portfolios.}
The CAPM \parencite{sharpe1964}, Roll's critique \parencite{roll1977} and its
correlation-aware refinements \parencite{prono2009}, fundamental indexing
\parencite{hsu2006}, covariance shrinkage \parencite{ledoitwolf2004},
tracking-error analysis \parencite{roll1992}, and the Markowitz optimization
baseline \parencite{markowitz1952} are taken up where they bear on the argument
(\Cref{sec:capm,sec:robust}).

\paragraph{Perturbed optimizers.}
The identity ``weights equal the gradient of an expected maximum'' is the
random-utility social-surplus construction \parencite{mcfadden1981, williams1977}
and, in machine learning, the perturbed-optimizer / differentiable-argmax
framework \parencite{berthet2020}. Gumbel perturbations give the
entropy-regularized softmax; our Gaussian, correlation-bearing perturbation gives
the correlation-aware regularizer of \Cref{thm:objective}.

\section{Discussion}
\label{sec:discussion}

The ability tilt is a mechanically simple idea with useful properties: long-only and feasible
without optimization, never inverting a covariance, reproducing a chosen benchmark
by construction and tilting only in response to correlation, and solving an
interpretable correlation-aware regularized program (\Cref{thm:objective}) whose
weights vary smoothly enough (\Cref{thm:smooth}) to bound turnover and license
cheap online updates and gradient-based tuning.

Open questions include the choice of reference correlation $C_{\mathrm{calib}}$
(diagonal versus market) and of $\phi$; the treatment of estimation error in the
abilities, a Thurstonian analogue of covariance shrinkage; and richer
perturbations, such as a temperature controlling concentration or using
$C_{\mathrm{tilt}}$ to express an explicit correlation view rather than the sample
estimate.

The tail-consistency result (\Cref{prop:tail}, \Cref{rem:anysim}) turns the
fat-tailed-performances direction into a concrete programme: drive the race with an
online \emph{downside} covariance (lower-partial-moment) estimator, which
\Cref{fig:tail} shows recovers almost all of the effect, and reserve the full
tail-dependent simulation for the residual. Estimating that downside covariance
robustly and online, and measuring whether the tail-aware tilt pays out of sample
on drawdown and expected shortfall rather than Sharpe, is the natural next study.

